\section{Data Splitting}

\subsection{Split Strategy}

The Corona NLP dataset provides a fixed official test file (\path{Corona_NLP_test.csv}, 3,798 samples) that is kept completely held-out and evaluated only once at the end of training. The official training file (41,157 samples) is further partitioned into a training subset and a validation subset.

We apply a \textbf{stratified random split} with a 90/10 ratio and a fixed random seed of 42, yielding:

\begin{equation}
    \underbrace{41{,}157}_{\text{official train}} \;\longrightarrow\;
    \underbrace{37{,}041}_{\text{train}} \;+\; \underbrace{4{,}116}_{\text{val}}
\end{equation}

Stratification ensures that the class distribution in the validation set mirrors that of the full training file, preventing any class from being under- or over-represented in the validation fold. Table~\ref{tab:split_summary} summarises the resulting split sizes.

\begin{table}[H]
    \centering
    \caption{Summary of dataset splits used in all experiments.}
    \label{tab:split_summary}
    \begin{tabularx}{\textwidth}{l r r l}
        \toprule
        \textbf{Split} & \textbf{Samples} & \textbf{Share (\%)} & \textbf{Purpose} \\
        \midrule
        Train      & 37,041 & 90.0 & Model parameter optimisation \\
        Validation & 4,116  & 10.0 & Checkpoint selection, early stopping \\
        Test        & 3,798  & ---  & Final one-time evaluation \\
        \midrule
        \textbf{Total (w/o test)} & \textbf{41,157} & \textbf{100.0} & \\
        \bottomrule
    \end{tabularx}
\end{table}

\subsection{Data Leakage Prevention}

The following measures are taken to prevent leakage between splits:

\begin{itemize}
    \item \textbf{Tokenizer is pretrained.} The DistilBERT tokenizer vocabulary is fully specified before any data is seen; there is no vocabulary fitting step that could leak validation or test token statistics into training.
    \item \textbf{Split before any data-dependent processing.} The train/val partition is performed first; all subsequent statistics (batch-level padding lengths, class counts for logging) are computed only on the training subset.
    \item \textbf{Validation is used only for checkpoint selection.} Validation macro-F1 is used to select the best model checkpoint and to trigger early stopping; no hyperparameter search loops over the validation set.
    \item \textbf{Test is evaluated exactly once.} The official test file is loaded only after training concludes and the best checkpoint is selected. No decisions are made based on test-set performance.
\end{itemize}

\subsection{Reproducibility}

All randomness is controlled by a single global seed (42), which is set for \texttt{random}, \texttt{numpy}, \texttt{torch}, and \texttt{torch.cuda} at the start of every run. The \texttt{sklearn.model\_selection.train\_test\_split} function with \texttt{stratify=labels} and \texttt{random\_state=42} produces identical splits across runs and machines.
